\title{Direct Preference Optimization: Your Language Model is Secretly a Reward Model}

\begin{abstract}
Large-scale unsupervised language models (LMs) acquire extensive world knowledge and develop certain reasoning abilities through training, yet exerting fine-grained control over their output remains challenging owing to the entirely unsupervised nature of their development. To address this limitation, common approaches collect human judgments on the comparative quality of different model outputs and subsequently adjust the LMs to favor these preferences, frequently utilizing reinforcement learning from human feedback (RLHF). Despite their effectiveness, RLHF methods introduce significant complexity and instability: they involve first training a reward model to mirror human preferences and then applying reinforcement learning to fine-tune the initial LM so as to maximize the learned reward, all while attempting to prevent the model from diverging excessively from its pretrained state. In this work, we propose a novel reward model parameterization for RLHF, which allows for direct derivation of the associated optimal policy in closed form. This advancement makes it possible to address the conventional RLHF objective using a straightforward classification loss. The proposed method, termed \textit{Direct Preference Optimization} (DPO), offers stability, strong performance, and low computational demands, removing the necessity for language model sampling during fine-tuning or for elaborate hyperparameter tuning. Experimental results indicate that DPO fine-tunes LMs to reflect human preferences at least as well as, or better than, alternative methods. Notably, DPO surpasses PPO-based RLHF in modulating generational sentiment and achieves comparable or superior outcomes in summarization and single-turn dialogue tasks, all while being considerably more straightforward to train and deploy.
\end{abstract}

\section{Introduction}
Large-scale unsupervised language models (LMs), when trained on massive corpora, demonstrate a range of unexpected capabilities~\citep{chowdhery2022palm, brown2020language, touvron2023llama,bubeck2023sparks}. These models, however, learn from human-generated data encompassing diverse intentions, objectives, and levels of expertise. Not all of these characteristics are suitable for replication; for instance, it is beneficial for an AI coding assistant to \textit{recognize} frequent programming errors so that it can rectify them. Yet, during code generation, we prefer that the model favor the (perhaps infrequent) exemplary programming practices embedded in its data. Analogously, it is useful for a language model to be \textit{informed} about a widespread misconception believed by half the population, but we decidedly do not want it to propagate this falsehood in half of the related queries. In summary, it is vital to distinguish and elicit the model’s \emph{intended responses and behavior} from the broad expanse of its \textit{learned knowledge and skills} for constructing AI systems that are safe, high-performing, and manageable \citep{ouyang2022training}. While current strategies generally rely on reinforcement learning (RL) to guide LMs toward human-aligned preferences, we demonstrate that the objective underlying these RL-based approaches can be met exactly using a straightforward binary cross-entropy loss, thereby streamlining the preference learning process.

At a conceptual level, prevalent techniques encode desirable model behaviors through carefully assembled datasets of human preference annotations, capturing the forms of conduct deemed safe and useful by humans. This stage of preference alignment follows an initial phase of extensive unsupervised pre-training on large text corpora. Although a naive method for preference learning is to apply supervised fine-tuning to human-provided examples of high-quality outputs, the most effective approaches employ reinforcement learning from human (or AI-generated) feedback (RLHF/RLAIF; \citep{christiano2017deep,bai2022constitutional}). RLHF involves training a reward model on a set of human preferences, which is subsequently used in an RL framework to adjust the language model’s policy, encouraging responses that achieve high reward while limiting divergence from the original pre-trained model. Despite the impressive performance of RLHF models, especially in dialogue and programming tasks, the RLHF pipeline is substantially more complicated than conventional supervised learning. It necessitates training several LMs and performing on-policy sampling throughout the learning loop, leading to considerable computational demands.

In this work, we introduce a method for optimizing language models to align with human preferences without depending on explicit reward models or reinforcement learning techniques. We present \textit{{Direct Preference Optimization} (DPO)}, an approach that implicitly optimizes for the same goal as conventional RLHF algorithms—that is, maximizing reward while enforcing a KL-divergence constraint—yet remains straightforward both to implement and train. Conceptually, the DPO mechanism boosts the log probability of preferred responses relative to those not favored, employing an adaptive, sample-specific importance weighting scheme that mitigates the issues of model collapse often observed when simply relying on a raw probability ratio objective. Similar to prior algorithms, DPO utilizes a formal preference modeling framework (for example, the Bradley-Terry model; \cite{bradley1952rankanalysis}) to gauge how closely a reward function matches observed preference judgments. The key distinction is that, unlike existing approaches—which train a reward model from preference losses and subsequently use it to train a policy—DPO reformulates the preference loss directly in terms of the policy via a variable transformation. Consequently, given a collection of human preference annotations over model outputs, DPO is able to train a policy using a binary cross entropy loss, thereby yielding an optimal policy corresponding to an implicitly learned reward function tailored to the preference dataset.

Our principal contribution is Direct Preference Optimization (DPO), an accessible algorithm for preference-based language model training that operates without reinforcement learning. Empirically, we demonstrate that DPO matches or exceeds the performance of established techniques, including RLHF methods based on PPO, for preference-driven learning across tasks such as sentiment control, summarization, and conversation, using language models with parameter counts as large as 6B.

Increasingly large self-supervised language models have demonstrated the ability to perform certain tasks without any labeled examples—zero-shot—or with minimal instruction via few-shot prompting \citep{radford2019language, gpt3, megatron, chowdhery2022palm}. Nonetheless, their effectiveness on specific tasks and their alignment with user goals can be substantially enhanced through fine-tuning on datasets comprised of instructions and responses authored by humans \citep{mishra-etal-2022-cross, sanh2022multitask, chung2022scaling, thoppilan2022lamda}. This process, known as `instruction-tuning', allows large language models to extrapolate to new instructions beyond those encountered during fine-tuning, thereby improving their overall utility \citep{chung2022scaling}. While instruction tuning has achieved considerable success, gathering \textit{relative} human preference data for response quality is generally less demanding than acquiring expert-labeled demonstrations. As a result, later approaches have focused on fine-tuning LLMs using datasets that capture human preferences, which has led to improved performance in domains such as translation \citep{kreutzer-etal-2018-reliability}, summarization \citep{stiennon2022learning, ziegler2020finetuning}, story generation \citep{ziegler2020finetuning}, and following instructions \citep{ouyang2022training, ramamurthy2023is}. These strategies typically involve an initial step where a neural reward function is trained to be consistent with the collected preferences, often modeled through frameworks such as the Bradley-Terry model \citep{bradley1952rankanalysis}. Subsequently, a language model is further adapted to optimize the reward signal, commonly utilizing reinforcement learning algorithms including REINFORCE \citep{williams1992reinforce}, proximal policy optimization (PPO; \cite{schulman2017proximal}), or related methods \citep{ramamurthy2023is}. Related research has also explored leveraging LLMs that have already been instruction-tuned with human feedback to synthesize new preference data targeting particular qualities such as safety or harmlessness \citep{bai2022constitutional}, relying on weak supervision—such as textual rubrics—to guide model-generated evaluations. Collectively, these developments represent a merging of research traditions: one focused on the application of reinforcement learning to train language models across diverse objectives \citep{Ranzato2015SequenceLT, paulus2018a, wu2018learning}, and another emphasizing methodologies for learning from human preference data \citep{christiano2017deep, kupcsik2018learning}. Notwithstanding the advantages of using relative human feedback, implementing reinforcement learning for fine-tuning large language models remains logistically complex; the present work introduces a principled method for preference optimization that avoids the need for RL.

Beyond language-specific applications, research on learning policies from preferences spans both bandit problems and reinforcement learning, with multiple methodologies explored. In the contextual bandit paradigm, where feedback consists of preferences or rankings over actions rather than explicit reward signals, the problem is framed as a contextual dueling bandit (CDB; \cite{yue2012karmed,dudik2015contextual}). Since absolute rewards are not observed, the analysis of CDBs introduces the concept of a \textit{von Neumann winner}: a policy guaranteed to have an expected win probability of at least 50\% when compared to any other policy \citep{dudik2015contextual}. However, in CDBs, preference feedback is typically provided in an online manner, whereas learning from human preferences often involves utilizing a static, pre-collected set of annotated action pairs \citep{yan2022human}. Along related lines, \textit{preference-based RL} (PbRL) methods rely on binary preference signals produced by an unobserved ‘scoring’ function instead of using explicit rewards \citep{BusaFekete2014,ruiz2023dueling}. Numerous PbRL algorithms exist, including those that leverage off-policy preference datasets, but these generally entail a two-step process: first learning the hidden scoring (reward) function, then optimizing the agent with respect to this estimate \citep{jain2013learning,BusaFekete2014,christiano2017deep,sadigh2017active,kupcsik2018learning}. In contrast, our work introduces a single-phase approach that directly trains a policy to conform to observed preferences.

\section{Preliminaries}\label{section:prelims}

We briefly recap the RLHF framework as implemented in \citeauthor{ziegler2020finetuning} and subsequent extensions (\citep{stiennon2022learning, bai2022training, ouyang2022training}). The process typically comprises three key steps: (1) supervised fine-tuning (SFT), (2) collecting preference data and fitting a reward model, and (3) reinforcement learning-based policy optimization.

\textbf{SFT}: The RLHF process generally initiates with the supervised fine-tuning of a large pre-trained language model on curated, high-quality examples specific to the target tasks (such as conversational agents, summarization, etc.), yielding an adapted model denoted as $\pi^\text{SFT}$.

\textbf{Reward Modelling Phase}: Next, the fine-tuned model is given prompts $x$ and generates answer pairs $(y_1, y_2)\sim \pi^\text{SFT}(y \mid x)$. These responses are then evaluated by human annotators, who select their preferred completion, expressed as $y_w\succ y_l \mid x$, where $y_w$ and $y_l$ indicate the selected and rejected answers from $(y_1, y_2)$, respectively. The preference annotation is modeled as being sampled from some latent reward function $r^*(y, x)$ that remains unobserved. Multiple techniques have been employed to represent such preferences, with the Bradley-Terry (BT) model \cite{bradley1952rankanalysis} being a prevalent option; more flexible models such as Plackett-Luce ranking systems \citep{plackett1975analysis, luce2012individual} are also compatible when multiple ranked alternatives are available. The BT formulation specifies that the human preference distribution $p^*$ takes the following form:

Given a static set of comparison data $\mathcal{D}=\bigl\{x^{(i)}, y_w^{(i)}, y_l^{(i)}\bigr\}_{i=1}^N$ drawn according to $p^*$, it is possible to represent the reward with a model $r_{\phi}(x, y)$, and infer the parameters by maximizing the likelihood. When the task is recast as binary classification, the following negative log-likelihood loss is minimized:

\begin{equation}\label{eq:reward_model}
    \mathcal{L}_R(r_{\phi}, \mathcal{D}) = -\mathbb{E}_{(x, y_w, y_l)\sim \mathcal{D}}\bigl[\log \sigma(r_{\phi}(x, y_w)- r_{\phi}(x, y_l))\bigr]
\end{equation}

where the logistic sigmoid $\sigma$ operates on the difference in predicted reward. Within language modeling frameworks, $r_{\phi}(x, y)$ is usually initialized from the SFT policy $\pi^\text{SFT}(y \mid x)$, extending it with a linear readout layer atop the final transformer block, so as to produce a scalar reward prediction for each input-output pair \cite{ziegler2020finetuning}. In an effort to reduce reward variance, some earlier works normalize reward outputs so that $\mathbb{E}_{x,y\sim \mathcal{D}}\left[r_\phi(x, y)\right] = 0$ for all $x$.

\textbf{RL Fine-Tuning Phase}: Once the reward function has been fit, it is used during reinforcement learning to supply feedback to the model. As outlined in prior research~\citep{jaques2017sequence, jaques2020human}, the following optimization is considered:

\begin{equation}\label{eq:RL}
\max_{\pi_{\theta}}  \mathbb{E}_{x\sim \mathcal{D}, y\sim \pi_{\theta}(y \mid x)}\bigl[r_{\phi}(x, y)\bigr] - \beta\mathbb{D}_{\textrm{KL}}\bigl[\pi_{\theta}(y\mid x)\mid \mid \pi_\text{ref}(y\mid x)\bigr],
\end{equation}

where $\beta$ tunes the regularization toward the reference distribution $\pi_\text{ref}$, corresponding in practice to the SFT model $\pi^\text{SFT}$. The policy $\pi_\theta$ is generally initialized from $\pi^\text{SFT}$ as well. Regularizing with the KL-divergence serves to anchor the fine-tuned model within the data distribution on which the reward predictions remain reliable, thus preserving sample diversity and guarding against mode collapse into a narrow set of high-reward outputs. Since the non-differentiable nature of discrete sequence generation prohibits end-to-end gradient optimization, reinforcement learning algorithms are typically employed to maximize this objective. The prevalent approach \citep{ziegler2020finetuning, stiennon2022learning, bai2022training, ouyang2022training} sets the reward as ${r(x, y) = r_{\phi}(x, y) -\beta (\log \pi_{\theta}(y\mid x) - \log \pi_\text{ref}(y\mid x))}$, with maximization performed using PPO \cite{schulman2017proximal}.

\section{Direct Preference Optimization}\label{sec:DPO}

Acknowledging the significant computational and practical hurdles in deploying reinforcement learning on massive language model fine-tuning tasks, we seek a more streamlined preference-driven policy optimization technique. Departing from the reward-then-RL paradigm of earlier RLHF strategies, our method utilizes a tailored reward model parameterization that enables closed-form computation of the corresponding optimal policy, thus obviating the need for iterative RL updates. The cornerstone of this method is an explicit mapping between reward functions and optimal policies, which we exploit to reformulate an objective in terms of policy parameters, effectively sidestepping the necessity of a standalone, separately-trained reward model—while still honoring established human preference models such as the Bradley-Terry formulation. Here, the policy network is crafted to encode both the conditional language generation and its latent preference-aware reward signal.

\textbf{Reformulating the DPO objective.} We commence from the standard RL objective provided in Eq.~\ref{eq:RL}, parameterized by a general reward function $r$. Prior literature~\citep{peters2007reinforcement, peng2019advantage, korbak2022reinforcement, go2023aligning} establishes that the solution maximizing the reward under a KL constraint, as in Eq.~\ref{eq:RL}, admits the following form:

\begin{equation}\label{eq:op_policy}
    \pi_r(y\mid x) = \frac{1}{Z(x)}\pi_\text{ref}(y\mid x)\exp\left(\frac{1}{\beta}r(x, y)\right),
\end{equation}

where the normalization constant $Z(x) =\sum_{y}\pi_\text{ref}(y\mid x)\exp\left(\frac{1}{\beta}r(x, y)\right)$ denotes the partition function. A detailed derivation can be found in Appendix \ref{app:derivation1}. Although employing an MLE-based approximation $r_{\phi}$ for the true reward $r^*$ is possible, accurately computing $Z(x)$ is generally computationally intensive~\citep{korbak2022reinforcement, go2023aligning}. This complicates practical applications of this policy representation. Nonetheless, we can algebraically manipulate Eq.~\ref{eq:op_policy} to re-express the reward as a function of the optimal policy $\pi_r$, the reference policy $\pi_\text{ref}$, and the (unknown) partition function $Z(\cdot)$. Taking the logarithm of Eq.~\ref{eq:op_policy} and rearranging gives:

\begin{equation}\label{eq:main_eq}
    r(x,y) =\beta \log \frac{\pi_r(y\mid x)}{\pi_\text{ref}(y\mid x)} + \beta \log Z(x).
\end{equation}

Applying this change of variables to the true reward $r^*$ and its associated optimal policy $\pi^*$, we note that the Bradley-Terry model is sensitive only to reward differences between completions; i.e., ${p^*(y_1 \succ y_2 \mid x) = \sigma(r^*(x, y_1) - r^*(x, y_2))}$. Substituting the reparameterization from Eq.~\ref{eq:main_eq} for $r^*(x,y)$ into the preference model (Eq.~\ref{eq:bradley-terry}) eliminates the partition function, enabling us to recast the human preference probability purely in terms of the optimal policy $\pi^*$ and the reference policy $\pi_\text{ref}$. Therefore, the optimal RLHF policy $\pi^*$ under the Bradley-Terry framework obeys:

\begin{equation}\label{eq:objective}
    p^*(y_1\succ y_2 \mid x)=\frac{1}{1 + \exp\left(\beta \log \frac{\pi^*(y_2\mid x)}{\pi_\text{ref}(y_2\mid x)} - \beta \log \frac{\pi^*(y_1\mid x)}{\pi_\text{ref}(y_1\mid x)}\right)}
\end{equation}

We provide a full derivation in Appendix~\ref{app:derivation2}. While Eq.~\ref{eq:objective} leverages the Bradley-Terry model, a parallel argument extends to the broader Plackett-Luce class~\citep{plackett1975analysis, luce2012individual}, as shown in Appendix~\ref{app:plackett_luce_models}.

Given this preference likelihood now depends on the policy itself rather than the reward function, we may directly optimize the maximum likelihood over a parameterized policy $\pi_\theta$. In a manner reminiscent of the approach for reward models (cf. Eq.~\ref{eq:reward_model}), the corresponding objective for the policy becomes:

\begin{equation}\label{eq:optimum_model}
    \mathcal{L}_\text{DPO}(\pi_{\theta}; \pi_\text{ref}) = -\mathbb{E}_{(x, y_w, y_l)\sim \mathcal{D}}\left[\log \sigma \left(\beta \log \frac{\pi_{\theta}(y_w\mid x)}{\pi_\text{ref}(y_w\mid x)} - \beta \log \frac{\pi_{\theta}(y_l\mid x)}{\pi_\

By employing this method, we learn an implicit reward function via an alternative parameterization, where the optimal policy corresponds directly to $\pi_\theta$. Additionally, as this approach can be viewed as optimizing a reparametrized Bradley-Terry model, it inherits several theoretical guarantees, including consistency when the preference data is sampled under appropriate conditions \cite{bong2022generalized}. We elaborate on DPO’s theoretical characteristics and connections to related methods in Section~\ref{sec:theory}.

\textbf{How does the DPO update operate?} To gain insight into the mechanics of DPO, it is beneficial to examine the gradient of the objective $\mathcal{L}_\text{DPO}$. The gradient with respect to the parameters $\theta$ is given by:

\begin{multline*}\label{eq:gradient}
    \nabla_\theta \mathcal{L}_\text{DPO}(\pi_\theta;\pi_\text{ref}) = \\ -\beta\mathbb{E}_{(x, y_w, y_l) \sim \mathcal{D}} \bigg[\underbrace{\sigma(\hat{r}_\theta(x, y_l) - \hat{r}_\theta (x, y_w))}_\text{emphasis when reward assignment is incorrect}\bigg[\underbrace{\nabla_\theta\log \pi(y_w \mid x)}_\text{boost preference for $y_w$} - \underbrace{\nabla_\theta\log\pi(y_l \mid x)}_\text{reduce preference for $y_l$}\bigg]\bigg],
\end{multline*}

with $\hat{r}_\theta(x, y) = \beta \log \frac{\pi_\theta(y \mid x)}{\pi_\text{ref}(y \mid x)}$ defining the reward implicitly assigned by the language model $\pi_\theta$ in comparison to the reference model $\pi_\text{ref}$ (further discussed in Section~\ref{sec:theory}). Conceptually, this gradient expression serves to enhance the probability of preferred responses $y_w$ while suppressing those of the less-preferred responses $y_l$. Notably, the magnitude of each update is modulated by the extent to which the implicit reward function $\hat{r}_\theta$ erroneously favors the dispreferred completion, modulated by $\beta$, thereby reflecting both the model’s error and the intensity of the KL constraint. Our experimental findings highlight that this weighting is critical: removing this factor leads to degenerate model behavior, as demonstrated in Appendix Table~\ref{tab:unlikelihood_generations}.

\textbf{DPO pipeline overview.} The typical DPO procedure involves the following steps: 1) For each prompt $x$, generate completion samples $y_1, y_2 \sim \pi_\text{ref}(\cdot \mid x)$, and use human preference annotations to assemble an offline dataset of preferences, denoted $\mathcal{D} = \{x^{(i)}, y_w^{(i)}, y_l^{(i)}\}_{i=1}^N$; 2) Train the language model $\pi_\theta$ to minimize the objective $\mathcal{L}_\text{DPO}$ with respect to the fixed $\pi_\text{ref}$, using the dataset $\mathcal{D}$ and the desired value of $\beta$.

In practical settings, there is often a preference to leverage existing publicly available preference datasets instead of generating new samples and collecting human feedback. Because these datasets are usually derived from $\pi^\text{SFT}$, it is standard to set $\pi_\text{ref} = \pi^\text{SFT}$ whenever such a model is accessible. Conversely, in situations where $\pi^\text{SFT}$ cannot be obtained, we set $\pi_\text{ref}$ by maximizing the likelihood of the preferred completions $(x, y_w)$, i.e., ${\pi_\text{ref} = \argmax_{\pi}\mathbb{E}_{x, y_w \sim \mathcal{D}}\left[\log \pi(y_w \mid x)\right]}$. This approach is intended to reduce the discrepancy between the unknown actual reference distribution and the chosen $\pi_\text{ref}$ employed in DPO. Additional specifics about implementation and selected hyperparameters are presented in Appendix~\ref{app:implementation}.

\section{Theoretical Analysis of DPO}

In this section, we provide further insight into the DPO approach, including theoretical justification and discussion of its strengths, particularly in relation to challenges present in actor-critic algorithms commonly used for RLHF (such as PPO~\cite{schulman2017proximal}).

\label{sec:theory}

\subsection{Your Language Model Is Secretly a Reward Model}

DPO achieves both reward learning and policy optimization in a unified maximum likelihood formulation, eliminating the need to explicitly fit a reward function or execute an RL procedure. The objective in Eq. \ref{eq:main_eq} is mathematically equivalent to the Bradley-Terry model, with a reward formulation of $r^*(x, y) = \beta \log\frac{\pi^*_\theta(y \mid x)}{\pi_\text{ref}(y \mid x)}$, and training the parametric model $\pi_{\theta}$ corresponds to optimizing the reward model objective (Eq. \ref{eq:reward_model}) under an appropriate variable transformation. In this section, we elaborate on the theoretical foundation for this reparameterization, demonstrate that it does not restrict the expressivity of the class of reward models that can be learned, and show that it facilitates precise recovery of the optimal policy. We begin by introducing an equivalence relation between reward functions.

\begin{definition}
Two reward functions $r(x, y)$ and $r'(x, y)$ are considered equivalent if and only if there exists a function $f$ such that $r(x, y) - r'(x, y) = f(x)$.
\end{definition}

It is straightforward to verify that this forms an equivalence relation, thereby dividing the collection of reward functions into distinct equivalence classes. This leads to the following two lemmas:

\begin{lemma}\label{lemma:same_prefrence} Within the context of the Plackett-Luce family—including, in particular, the Bradley-Terry model—any pair of reward functions from the same equivalence class yield identical distributions over preferences.
\end{lemma}

\begin{lemma}\label{lemma:same_policy}
    Within the framework of constrained RL, two reward functions belonging to the same equivalence class will result in the same optimal policy.
\end{lemma}

Given the straightforward nature of their proofs, we provide these in Appendix~\ref{app:lemma1}. The first lemma addresses the classic under-determination present in Plackett-Luce models \cite{plackett1975analysis}. As a result of this ambiguity, it is often necessary to introduce additional constraints for identifiability when deriving MLE estimates based on Eq.~\ref{eq:reward_model} \cite{bong2022generalized}. The second lemma indicates that all members of an equivalence class determine the same optimal policy, which implies that for the purpose of optimizing our objective, it suffices to recover any representative reward function from the equivalence class corresponding to the optimum. The subsequent theorem, proved in Appendix~\ref{app:thm1}, formalizes this observation:

\begin{theorem}\label{thm:main}
    Assuming mild regularity conditions, every reward equivalence class compatible with the Plackett-Luce (and in particular the Bradley-Terry) family can be expressed via the reparameterization ${r(x, y) = \beta \log \frac{\pi(y\mid x)}{\pi_\text{ref}(y\mid x)}}$ for some policy $\pi(y\mid x)$ and a specified reference policy $\pi_\text{ref}(y \mid x)$.
\end{theorem}

\begin{sproof}
    Given a reward function $r(x, y)$, let its corresponding optimal policy be $\pi_r(y \mid x)$, as prescribed by Eq.~\ref{eq:op_policy}. We aim to demonstrate that any reward function within the equivalence class of $r$ admits the reparameterized form stated above. Define the projection
\begin{equation}
    f(r; \pi_\text{ref}, \beta)(x, y) = r(x, y) - \beta\log\sum_{y}\pi_\text{ref}(y\mid x)\exp\left(\frac{1}{\beta}r(x, y)\right)
\end{equation}
The operator $f$ applies normalization by subtracting the log-partition function associated with $\pi_r$. Since the normalization is solely a function of $x$, it follows that $f(r; \pi_\text{ref}, \beta)(x, y)$ resides in the equivalence class of $r(x, y)$. Further, if we substitute $r$ with the right-hand side of Eq.~\ref{eq:main_eq}—a form valid for any $r$—we obtain $f(r; \pi_\text{ref}, \beta)(x, y) = \beta \log \frac{\pi_r(y\mid x)}{\pi_\text{ref}(y\mid x)}$. This confirms that the projection $f$ recovers a canonical representative of the equivalence class, ensuring that our choice of reparameterization for the reward function model does not lose generality.
\end{sproof}

An alternative interpretation of Theorem~\ref{thm:main} is that it pinpoints the precise reward function chosen by the DPO reparameterization from each equivalence class, specifically the one that satisfies:

\begin{equation}\label{eq:lag_p}
     \sum_{y}\underbrace{\pi_\text{ref}(y\mid x)\exp\left(\frac{1}{\beta}r(x, y)\right)}_{=\pi(y\mid x)\text{, using Thm.~\ref{thm:main} reparam.}} = 1,
\end{equation}

In other words, this ensures that $\pi(y\mid x)$ constitutes a legitimate probability distribution, meaning its values are non-negative and they collectively sum to one. Building on Eq.~\ref{eq:op_policy}, we observe that Eq.~\ref{eq:lag_p} actually defines the partition function associated with the optimal policy that is derived from the reward function $r(x, y)$. A central insight behind the DPO algorithm is the realization that, by enforcing suitable constraints on the inherently under-specified Plackett-Luce (and especially the Bradley-Terry) class of preference models, it is possible to both retain the expressive power for representing various reward models and guarantee that the optimal policy from Eq.~\ref{eq:op_policy} is tractable in closed-form for any given prompt $x$.

\subsection{Instability of Actor-Critic Algorithms} Our analytical framework also enables us to elucidate why standard actor-critic approaches, such as PPO, can become unstable in RLHF applications. Specifically, we focus on the reinforcement learning fine-tuning phase described in Section~\ref{section:prelims}, consistent with the standard RLHF pipeline. This setup aligns closely with the control as inference perspective \cite{levine2018reinforcement} when applied to the constrained RL problem set out in \ref{eq:RL}. Assuming a parameterized distribution $\pi_{\theta}(y\mid x)$, the aim is to minimize $\mathbb{D}_{\text{KL}}[\pi_{\theta}(y|x) \mid \mid \pi^*(y\mid x)]$, where $\pi^*$ represents the optimal policy induced by the reward $r_{\phi}(y, x)$, as specified by Eq.~\ref{eq:optimum_model}. Through further manipulation, we arrive at the following optimization criterion:

\begin{equation}\label{eq:AC}
    \max_{\pi_{\theta}}\mathbb{E}_{\pi_{\theta}(y\mid x)}\bigg[\underbrace{r_{\phi}(x, y) -\beta\log\sum_{y}\pi_\text{ref}(y\mid x)\exp\left(\frac{1}{\beta}r_{\phi}(x, y)\right)}_{f(r_{\phi}, \pi_\text{ref}, \beta)} - \underbrace{\beta\log\frac{\pi_{\theta}(y\mid x)}{\pi_\text{ref}(y\mid x)}}_{\text{KL}}\bigg]
\end{equation}

The objective considered here is identical to that optimized in earlier research 
\citep{ziegler2020finetuning, stiennon2022learning, bai2022training, ouyang2022training}, employing a DPO-equivalent reward formulated for the reward class $r_{\phi}$. Within this framework, the normalization component within $f(r_{\phi}, \pi_\text{ref}, \beta)$ can be viewed as representing the soft value function of the reference policy $\pi_\text{ref}$. Although this normalization term is inconsequential to the optimal solution itself, omitting it can result in a policy gradient with elevated variance, potentially destabilizing the learning process. One way to address the normalization term involves the use of a learned value function, though this approach also presents optimization challenges. An alternative strategy adopted by prior studies is to normalize the rewards using a baseline derived from a human completion, which effectively constitutes a single-sample Monte-Carlo approximation of the normalizer. In comparison, the DPO reparameterization produces a reward function that obviates the need for any such baseline.

\section{Experiments}
This section provides an empirical assessment of the effectiveness of DPO for training policies directly from preference data. We first examine DPO within a controlled text-generation environment to investigate its efficiency in balancing reward maximization against KL-divergence minimization from the reference policy, comparing it to prevalent preference optimization techniques like PPO. Subsequently, DPO is tested on larger language models and more complex RLHF challenges, such as summarization and dialogue tasks. Our results indicate that, with minimal hyperparameter tuning, DPO often matches or outperforms strong baselines including RLHF using PPO and even surpasses approaches based on selecting the highest-reward trajectory out of $N$ samples via a learned reward model. Prior to discussing these findings in detail, we outline the experimental protocol; further methodological specifics can be found in Appendix~\ref{app:exp_details}.

\textbf{Tasks.} We conduct experiments across three distinct open-ended text generation tasks. Across these tasks, each algorithm is trained to optimize a policy based on a dataset of preferences, $\mathcal{D}=\bigl\{x^{(i)}, y_w^{(i)}, y_l^{(i)}\bigr\}_{i=1}^N$. In the \textbf{controlled sentiment generation} task, $x$ represents an initial segment of a movie review sourced from the IMDb dataset \cite{maas-EtAl:2011:ACL-HLT2011}; the goal for the policy is to produce a continuation $y$ with explicitly positive sentiment. For reliable benchmarking, we synthetically construct preference pairs for generated texts using a pre-trained sentiment classifier such that $p(\text{positive}\mid x,y_w)>p(\text{positive}\mid x,y_l)$. For supervised fine-tuning (SFT), GPT-2-large is fine-tuned until convergence on the training split of the IMDb reviews (see App~\ref{app:sentiment_details} for further information). For the \textbf{summarization} task, $x$ is a Reddit forum post, and the policy is expected to generate a summary $y$ highlighting the core points. In line with existing literature, we employ the Reddit TL;DR summarization dataset \citep{volske-etal-2017-tl} along with human-labeled preference data curated by \citeauthor{stiennon2022learning}. The SFT baseline is a model fine-tuned on forum post summaries written by humans,\footnote{\url{https://huggingface.co/CarperAI/openai_summarize_tldr_sft}} and reinforcement learning from human feedback is implemented using the TRLX framework \citep{leandro_von_werra_2023_7790115}. The human preference labels are based on outputs generated from a different, albeit similarly-trained, SFT model as per \citeauthor{stiennon2022learning}. The third setting, \textbf{single-turn dialogue}, takes $x$ as a human prompt—ranging from science questions to personal advice requests—and requires the policy to craft an appropriate and constructive reply $y$. We utilize the Anthropic Helpful and Harmless dialogue dataset \citep{bai2022training}, encompassing 170k interactions between people and a machine assistant. Each exchange is followed by two assistant-generated replies from a large, undisclosed model, and is annotated with the human’s preferred completion. Since an SFT checkpoint is not publicly provided for this corpus, we create one by fine-tuning an existing language model solely on the preferred completions.

\textbf{Evaluation.} We adopt two main strategies for evaluation in our experiments. To assess each algorithm's ability to optimize the constrained reward maximization objective, we analyze their respective frontiers of reward versus KL-divergence from the reference policy within the controlled sentiment generation environment; this is feasible given our access to the true reward function (implemented via a sentiment classifier). However, in practical deployment scenarios where the ground truth reward function is unavailable, we measure algorithm performance in terms of \textit{win rate} against a baseline policy. Here, GPT-4 serves as a stand-in for human evaluation when judging the quality of summaries and the helpfulness of single-turn dialogue responses, depending on the task. For the summarization task, our baseline consists of reference summaries from the test set; in the case of dialogue, the baseline is the preferred response from the test dataset. Prior research indicates that language models can serve as more effective automatic evaluators compared to existing metrics \citep{Chen2023ExploringTU}. Nonetheless, to support the validity of using GPT-4 for our assessments, we conduct a human subject study as described in Sec.~\ref{sec:human-judgments}. Our findings show that GPT-4's assessments are highly correlated with those of humans, with human-to-GPT-4 agreement rates typically matching or exceeding agreement between individual human annotators.

\textbf{Methods.} Alongside DPO, we benchmark several established techniques for training language models to reflect human preferences. For baseline prompting, we use zero-shot prompts with \textbf{GPT-J} \citep{gpt-j} for summarization, and two-shot prompts with \textbf{Pythia-2.8B} \citep{biderman2023pythia} for the dialogue task. We further evaluate the \textbf{SFT} model as well as \textbf{Preferred-FT}, which is created through supervised fine-tuning on the preferred completion $y_w$, drawing from either the SFT model (in sentiment control and summarization tasks) or a general language model (in single-turn dialogue tasks). We also include the \textbf{Unlikelihood} method~\citep{welleck2019neural}, which seeks to simultaneously maximize the likelihood of $y_w$ and suppress the likelihood of $y_l$; an optional scaling parameter $\alpha\in[0,1]$ is applied to the unlikelihood loss. We compare these to \textbf{PPO} \citep{schulman2017proximal}, using a reward function learned from preference annotations, and \textbf{PPO-GT}, which is provided with access to the true reward in the controlled sentiment setting. Within our sentiment experiments, we utilize both an out-of-the-box PPO-GT implementation \cite{leandro_von_werra_2023_7790115} and a customized variant that applies reward normalization and further optimizes hyperparameters to boost performance (these customizations are similarly applied to PPO with learned rewards). As another baseline, we implement \textbf{Best of $N$}: here, we sample $N$ candidate responses from the SFT model (or from Preferred-FT in the dialogue scenario), and return the response scoring highest according to a reward model trained on preference data. While this approach can achieve high performance, it is not practical for even moderate $N$, as it necessitates generating $N$ completions per query during testing.

\subsection{How well can DPO optimize the RLHF objective?}

In RLHF frameworks, the KL-constrained reward maximization objective is designed to maximize the reward signal while ensuring that the learned policy does not diverge excessively from the reference policy. Consequently, meaningful comparisons between different algorithms require consideration of both the achieved reward and the KL divergence; attaining marginally higher rewards at the cost of substantially larger KL divergence is often undesirable. Figure~\ref{fig:frontier-tldr-main} displays the reward-KL tradeoff frontier across a range of algorithms for the sentiment task. For each algorithm, we conduct several independent runs, each employing distinct values for the policy conservativeness hyperparameter (specifically, target KL $\in\{3,6,9,12\}$ for PPO, $\beta \in \{0.05,0.1,1,5\}$, and $\alpha\in\{0.05,0.1,0.5,1\}$ for unlikelihood, and different random seeds for preferred-FT). In total, this experimentation covers 22 separate training runs. At intervals of 100 training steps, up until convergence, we assess every policy using a shared set of test prompts, recording both the mean reward under the true reward function and the average sequence-level KL\footnote{In other words, the cumulative sum of KL-divergences computed at each timestep.} with respect to the reference policy, denoted as $\text{KL}\left(\pi\mid \mid \pi_\text{ref}\right)$. Our observations indicate that DPO defines a significantly more favorable frontier, attaining superior reward outcomes without incurring high KL values. This finding is especially striking for several reasons. Most notably, even though both DPO and PPO target the same objective, DPO achieves greater efficiency; the reward/KL frontier for DPO is strictly superior to that of PPO. Moreover, DPO outperforms PPO's frontier, \emph{even under conditions where PPO is provided with access to the ground truth reward} (PPO-GT).

\subsection{Can DPO scale to real preference datasets?}
\label{sec:dpo-real-datasets}
We proceed to assess how well DPO fine-tunes models for tasks in summarization and single-turn dialogue. In the domain of summarization, commonly used automatic metrics like ROUGE often do not align closely with human judgment~\citep{stiennon2022learning}, and prior research has demonstrated that PPO-based fine-tuning guided by human preferences yields more useful summaries. To compare various approaches, we generate model outputs using the test partition of the TL;DR summarization dataset and calculate the mean win rate against ground-truth completions found in the test data. For each approach, completions are generated over a temperature range from 0.0 to 1.0, and the corresponding win rates appear in Figure~\ref{fig:frontier-tldr-main} (right). DPO, PPO, and Preferred-FT are all fine-tuned from the same GPT-J SFT model\footnote{\url{https://huggingface.co/CarperAI/openai_summarize_tldr_sft}}. Our experiments reveal that DPO achieves a win rate near 61\% at temperature 0.0, surpassing PPO, which achieves about 57\% at its best-performing temperature (0.0). Additionally, DPO attains a greater peak win rate when compared with the best of $N$ baseline. It is worth highlighting that DPO’s $\beta$ hyperparameter was not extensively optimized in these experiments, so actual performance might be understated. Furthermore, DPO exhibits significantly more stability across sampling temperatures relative to PPO, whose performance declines to that of the original GPT-J model as temperature increases. Preferred-FT, in contrast, does not show notable gains over the base SFT model. A direct human preference evaluation of DPO versus PPO, detailed in Section~\ref{sec:human-judgments}, shows that samples generated by DPO at temperature 0.25 are favored 58\% of the time when compared to PPO samples at the same temperature.

For the single-turn dialogue task, we assess the various approaches using the subset of the Anthropic HH dataset's test split \citep{bai2022training} that features a single exchange between human and assistant. For evaluation with GPT-4, we utilize the preferred completions from the test set as references, calculating the win rate across methods accordingly. Due to the absence of an established SFT baseline, we initiate experiments with a pre-trained Pythia-2.8B model. We employ Preferred-FT to develop a reference model aligned with the distribution of selected completions, and subsequently train using DPO. Additionally, we benchmark against both the strongest sample among 128 Preferred-FT completions (noting that the Best of $N$ strategy levels off at 128 completions, as detailed in Appendix Figure~\ref{fig:best-of-n}) and a 2-shot prompt version of the original Pythia-2.8B model, observing that DPO matches or surpasses the top-performing sampling temperatures of each method. We further include an RLHF model trained via PPO on the Anthropic HH dataset \footnote{\url{https://huggingface.co/reciprocate/ppo_hh_pythia-6B}}, referenced from a widely recognized implementation \footnote{\url{https://github.com/CarperAI/trlx/tree/main/examples/hh}}, but are unable to identify any prompt or sampling temperature configuration that exceeds the performance of the base Pythia-2.8B. Drawing from insights on the TL;DR task and acknowledging both Best of 128 and PPO share the same reward function, we interpret Best of 128 as an approximate stand-in for PPO results. In summary, DPO stands out as the only method among those evaluated that efficiently outperforms the preferred completions from the Anthropic HH dataset, offering results comparable to or better than the computationally intensive Best of 128 baseline. Lastly, Figure~\ref{fig:dialogue-main} illustrates DPO's rapid convergence toward optimal performance.

\subsection{Generalization to a new input distribution}

To more thoroughly assess how PPO and DPO handle shifts in input distribution, we evaluate the policies trained in our Reddit TL;DR summarization setting on a distinct domain: the test partition of the CNN/DailyMail news dataset \citep{nallapati-etal-2016-abstractive}. We maintain consistency by applying the same optimal sampling temperatures identified for TL;DR (0 and 0.25). Table~\ref{tab:ood} summarizes the performance results. For evaluation, we compute the win rate according to GPT-4 relative to the reference summaries, utilizing the same GPT-4 (C) prompt as before but with the term ``forum post'' substituted by ``news article''. On this novel distribution, DPO continues to significantly surpass PPO. These findings provide preliminary support for the notion that DPO models exhibit generalization capabilities comparable to those of PPO, despite not leveraging the additional unlabeled Reddit TL;DR prompts available to PPO.

\subsection{Validating GPT-4 judgments with human judgments}
\label{sec:human-judgments}
A human evaluation study is performed to assess the trustworthiness of GPT-4’s comparisons, using outputs from the TL;DR summarization task and two separate prompting strategies for GPT-4. The \textbf{GPT-4 (S)} (simple) prompt requests an assessment on which summary captures the core information of the post more effectively. By contrast, the \textbf{GPT-4 (C)} (concise) prompt asks not only for informativeness but also conciseness, motivated by our observation that the simple prompt tends to make GPT-4 favor verbose and repetitive summaries compared to human preferences. The exact prompt formulations are available in Appendix~\ref{app:prompts}. Our evaluation includes three model variants—DPO (temp. 0.25) as the top performer, PPO (temp. 1.0) as the lowest, and SFT (temp. 0.25) with intermediate quality—to ensure a representative range of sample outputs. All variants are benchmarked against PPO using greedy decoding, the strongest baseline among PPO temperatures. Across both prompting conditions, we observe that GPT-4’s choices align with human ratings approximately as consistently as the inter-human agreement rates, indicating that GPT-4 is a reasonable substitute for direct human assessment (although, because of the small human rater pool, multiple independent annotations are collected only for DPO and PPO-1). Notably, the \textbf{GPT-4 (C)} prompt tends to yield win rates that better mirror human preferences, prompting us to use it for the principal results reported in Section~\ref{sec:dpo-real-datasets}. Further methodological information about the human rating process, including interface screenshots and rater lists, is presented in Appendix~\ref{app:human-study}.

\section{Discussion}
Preference-based learning presents a robust and scalable approach for developing proficient and aligned language models. In this work, we introduced DPO, a straightforward methodology for learning from preferences without the necessity of reinforcement learning. Unlike traditional methods that reshape preference learning into an RL framework to utilize standard RL techniques, DPO establishes a correspondence between language model policies and reward functions. This allows language models to be trained to reflect human preferences \textit{directly}, employing a simple cross-entropy objective and bypassing both reinforcement learning and any loss of generality. Remarkably, DPO achieves performance on par with or superior to conventional RLHF approaches such as those utilizing PPO, and does so with minimal hyperparameter tuning. As a result, DPO significantly lowers the entry barriers for training language models informed by human feedback.

\textbf{Limitations \& Future Work.} Our findings open up numerous directions for further investigation. One key question is how DPO-trained policies handle out-of-distribution generalization relative to models trained on explicit reward signals; preliminary results indicate comparable generalization to PPO-trained models, yet extensive study is warranted. Further, an open issue is whether leveraging self-labeling using the DPO policy allows for effective exploitation of unlabeled prompts. Another research avenue pertains to understanding how reward over-optimization appears in direct preference optimization scenarios, and whether the marginal drop in performance in Figure~\ref{fig:dialogue-main}-right exemplifies this phenomenon. Additionally, while our analysis includes models with up to 6B parameters, expanding DPO to accommodate much larger, state-of-the-art architectures is a promising area for future exploration. In terms of evaluation, our work reveals that GPT-4-derived win rates are sensitive to prompt formulation; subsequent studies could investigate optimal methodologies for eliciting reliable assessments from automated evaluators. Lastly, DPO offers a broad array of prospective applications, not limited to human preference-aligned language models but also encompassing generative modeling across different domains.